\chapter{Αναφορά τοπικού API}
\label{app:api}

Ο ενσωματωμένος εξυπηρετητής ακούει στη θύρα 8080 της συσκευής. Κάθε αίτημα, εκτός από το \code{/api/health}, απαιτεί το κλειδί της συσκευής στην επικεφαλίδα \code{X-Api-Key} ή \code{Authorization: Bearer}. Σώματα και απαντήσεις είναι JSON, εκτός αν σημειώνεται διαφορετικά.

\section{Endpoints}

\begin{longtable}{@{}p{4.6cm}p{1.6cm}p{8cm}@{}}
  \caption{Endpoints του τοπικού API.}\label{tab:api-full}\\
  \toprule
  Διαδρομή & Μέθοδος & Λειτουργία \\
  \midrule
  \endfirsthead
  \toprule
  Διαδρομή & Μέθοδος & Λειτουργία \\
  \midrule
  \endhead
  \code{/api/health} & GET & Κατάσταση και χρόνος λειτουργίας (χωρίς κλειδί) \\
  \code{/api/status} & GET & Σημεία, ευθυγράμμιση, συνεδρία \\
  \code{/api/logs} & GET & Πρόσφατα αιτήματα (κυκλικός buffer) \\
  \code{/api/coordinates} & GET, POST & Λίστα / προσθήκη σημείων \\
  \code{/api/coordinates/{id}} & DELETE & Διαγραφή ενός σημείου \\
  \code{/api/clear} & POST & Διαγραφή όλων των σημείων \\
  \code{/api/undo} & POST & Αναίρεση της τελευταίας αλλαγής σημείων \\
  \code{/api/layer} & POST & Φίλτρο ορατού στρώματος \\
  \code{/api/project} & POST & Φόρτωση έργου ευθυγράμμισης \\
  \code{/api/project/reset} & POST & Επαναφορά της ευθυγράμμισης του έργου \\
  \code{/api/calibrate/start} & POST & Έναρξη βαθμονόμησης με σημεία \\
  \code{/api/calibrate} & POST & Αποστολή πραγματικών συντεταγμένων σημείου \\
  \code{/api/calibrate/status} & GET & Κατάσταση βαθμονόμησης και υπόλοιπο RMS \\
  \code{/api/calibrate/reset} & POST & Επαναφορά βαθμονόμησης \\
  \code{/api/gps-origin} & POST & Ορισμός γεωγραφικής αρχής \\
  \code{/api/anchors} & GET & Κατάσταση του anchor αναφοράς \\
  \code{/api/session} & GET, POST & Εξαγωγή / εισαγωγή συνεδρίας \\
  \code{/api/session/new} & POST & Νέα ονομασμένη συνεδρία \\
  \code{/api/model} & POST & Μεταφόρτωση μοντέλου (OBJ, DXF ή zip) \\
  \code{/api/models} & GET & Λίστα μοντέλων στη συσκευή \\
  \code{/api/fieldnotes} & GET & Σημειώσεις πεδίου \\
  \code{/api/fieldnotes/photo/{file}} & GET & Φωτογραφία σημείωσης πεδίου (JPEG) \\
  \code{/api/reports}, \code{/api/reports/{file}} & GET & Αναφορές διαδρομών επιθεώρησης \\
  \code{/api/fieldtest}, \code{/api/fieldtest/{file}} & GET & Αρχεία CSV της δοκιμής anchors \\
  \bottomrule
\end{longtable}

\review{Ο πίνακας είναι από το README και τη λίστα διαδρομών του \code{LocalAPIServer.cs}. Έλεγξε τις μεθόδους του \code{/api/model} (μεταφόρτωση, διαγραφή) και του \code{/api/project} στον κώδικα, και αν υπάρχουν endpoints που λείπουν.}

\section{Κωδικοί απάντησης}

\begin{table}[htbp]
  \centering\small
  \caption{Κωδικοί σφάλματος.}
  \label{tab:api-errors}
  \begin{tabularx}{\textwidth}{@{}lX@{}}
    \toprule
    Κωδικός & Πότε \\
    \midrule
    400 & Κενό ή μη έγκυρο JSON, μη έγκυρες τιμές (NaN, εκτός ορίων) \\
    401 & Λείπει ή είναι λάθος το κλειδί \\
    405 & Γνωστή διαδρομή με λάθος μέθοδο \\
    408 & Το σώμα δεν διαβάστηκε εγκαίρως \\
    413 & Σώμα JSON πάνω από 512\,KB ή μοντέλο πάνω από 300\,MB \\
    429 & Υπέρβαση ορίου ρυθμού (60/λεπτό χωρίς κλειδί, 600/λεπτό με κλειδί) ή αποκλεισμός μετά από 10 λάθος κλειδιά· επικεφαλίδα \code{Retry-After} \\
    503 & Η εφαρμογή δεν απάντησε εγκαίρως (π.χ. συσκευή σε αναμονή) ή η ουρά εντολών είναι γεμάτη· επικεφαλίδα \code{Retry-After} \\
    \bottomrule
  \end{tabularx}
\end{table}

\section{Παράδειγμα: αποστολή σημείων}

\begin{lstlisting}[language=bash,caption={Αποστολή σημείων με cURL.}]
curl -X POST http://<quest-ip>:8080/api/coordinates \
     -H "Content-Type: application/json" \
     -H "X-Api-Key: <key>" \
     -d '{"coordinates":[{"id":"K5","x":12.4,"y":1.2,"z":-3.8,
          "label":"K5","color":"#FF0000","scale":0.1,
          "markerType":"pin","layer":"hazards"}],
          "clearExisting":false,"removeMissing":false}'
\end{lstlisting}

Τα πεδία \code{x}, \code{y}, \code{z} είναι σε μέτρα, με το \code{y} προς τα πάνω. Το \code{clearExisting} διαγράφει όλα τα σημεία πριν από την εισαγωγή, και το \code{removeMissing} διαγράφει όσα λείπουν από το αίτημα (συγχρονισμός). Αντί για \code{x/y/z}, ένα σημείο μπορεί να δοθεί με \code{latitude}, \code{longitude}, \code{altitude} και \code{"crs":"gps"}, αφού οριστεί πρώτα αρχή με \code{/api/gps-origin}.
